% ===== wrap-up =====
\begin{frame}[fragile]{11주차 정리}
  \begin{enumerate}
    \item \kb{11.1 --- 은닉 상태는 ``감쇠하는 요약''} --- 같은 세
      행렬($W_{xh}, W_{hh}, W_{hy}$)을 $T$번 순서대로 돌려
      파라미터를 길이와 무관하게 일정하게 만들고, $w_{hh}$가
      ``기억의 길이''를 조절 --- \textbf{파라미터 공유 = CNN의
      시간 축 버전}.
    \item \kb{11.2 --- BPTT는 (A)+(B)의 재귀} --- 본 시점 오차 $+$
      다음 시점에서 온 그래디언트를 더하고, 공유 파라미터는
      \verb|+=|로 누적 --- ``abcabc…'' 8차원 문자모델은 LLM
      ``다음 토큰 예측''의 \textbf{가장 축소된 형태}(12.09
      $\to$ 0.00086, 14,000배).
    \item \kb{11.3 --- 곱은 지수, 덧셈은 통로} --- 그래디언트
      곱적 $\prod tanh'\, w_{hh}$가 고유값 1을 기준으로
      지수 소실·폭발; LSTM/GRU는 \textbf{곱(게이트)=조절,
      덧셈(셀 스킵)=통로}로 소실을 \textbf{완화} --- 단,
      \textbf{순차성이라는 근본 한계는 남고}, 이것이 Week 12
      어텐션의 계기.
  \end{enumerate}
  \vskip1em
  \begin{block}{다음 주 --- Week 12. 어텐션과 트랜스포머}
    RNN이 ``과거를 순차적으로 압축해서 기억''했다면, Attention은
    ``\textbf{필요할 때마다 과거 전체를 직접 다시 들여다본다}'' ---
    11.1의 인코더-디코더 한계(긴 입력을 하나의 벡터로 압축하면
    정보 손실)를 보완하러 등장하며, 11.3의 ``완화이지 해결이
    아니다''에 대한 \textbf{구조적} 답이다.
  \end{block}
\end{frame}
